\ssr{ВВЕДЕНИЕ}

HTTP-сервер для отдачи статических ресурсов является основным компонентом сетевой инфраструктуры. Он отвечает за выдачу неизменяемых ресурсов (HTML, CSS, JavaScript, изображения, файлы). Производительность такого сервера напрямую влияет на время отклика и суммарную нагрузку на вычислительные ресурсы.

Для обслуживания большого количества активных соединений серверу необходимо эффективно работать с сетевыми соединениями. Один из классических подходов~---~неблокирующий ввод-вывод и мультиплексирование по готовности дескрипторов~---~позволяет обслуживать множество соединений в рамках одного процесса.

Цель работы — разработать HTTP-сервер для отдачи статического содержимого с использованием мультиплексирования соединений на основе системного вызова select() и предварительно созданных процессов-обработчиков (worker’ов), обеспечивающих параллельную обработку соединений.

Для достижения поставленной цели необходимо решить следующие задачи:
\begin{itemize}
\item выполнить анализ предметной области: основные принципы устройства и работы протокола HTTP, механизм мультиплексирования с использованием системного вызова select() и подход к организации параллельной обработки соединений;
\item разработаны алгоритмы обработки HTTP-запросов и обслуживания множества соединений одним процессом-обработчиком;
\item выбрать средства реализации и реализовать HTTP-сервер;
\item провести исследование и определить: максимальное количество обслуживаемых сетевых соединений, зависимость скорости передачи данных по каждому сетевому соединению, совокупной средней скорости передачи и среднего времени обработки запроса.
\end{itemize}